```latex
\documentclass[11pt]{amsart}

\usepackage{amsmath,amssymb,amsthm}
\usepackage{geometry}
\usepackage{hyperref}
\geometry{margin=1.1in}

\theoremstyle{plain}
\newtheorem{theorem}{Theorem}[section]
\newtheorem{lemma}[theorem]{Lemma}
\newtheorem{proposition}[theorem]{Proposition}
\newtheorem{corollary}[theorem]{Corollary}

\theoremstyle{definition}
\newtheorem{definition}[theorem]{Definition}
\newtheorem{convention}[theorem]{Convention}
\newtheorem{hypothesis}[theorem]{Hypothesis}

\theoremstyle{remark}
\newtheorem{remark}[theorem]{Remark}

\newcommand{\Dp}{\mathcal{D}'}
\newcommand{\R}{\mathbb{R}}
\newcommand{\C}{\mathbb{C}}
\newcommand{\Prob}{\mathcal{P}}
\newcommand{\Lx}{\mathcal{L}_X}

\begin{document}

\title{Smooth solutions of the cohomological equation\\ force zero topological entropy}

\author{}
\date{}

\begin{abstract}
Let $M$ be a compact connected smooth manifold and let $X$ be a $C^\infty$
vector field on $M$ with flow $(g_t)_{t\in\R}$. We consider the following
regularity hypothesis on the cohomological equation: whenever
$f\in C^\infty(M)$ and $u\in\Dp(M)$ satisfy $\Lx u=f$ in the sense of
distributions, the distribution $u$ is smooth. We prove that this
hypothesis forces the flow to have zero topological entropy (where, as
is standard, the topological entropy of a flow is defined to be the
topological entropy of its time-one map $g_1$). Along the way we
establish strong structural consequences of the hypothesis which are of
independent interest: the flow is minimal, it is uniquely ergodic, its
unique invariant probability has a $C^\infty$ density which is strictly
positive everywhere, the space of invariant distributions is one
dimensional, and the topological entropy of $g_1$ equals the metric
entropy of the unique invariant measure. The final step of the proof is
the classical theorem of A.\,G.~Kushnirenko, according to which a smooth
flow with positive topological entropy admits non-smooth distributional
solutions of the cohomological equation.
\end{abstract}

\maketitle

\section{Introduction and statement of the results}

Let $M$ be a compact connected smooth ($C^\infty$) manifold without
boundary, let $X$ be a $C^\infty$ vector field on $M$, and let
$(g_t)_{t\in\R}$ denote the (complete, since $M$ is compact) flow
generated by $X$. Following the standard convention---used for instance
in the recent literature on entropies of smooth flows
\cite{LWW}---\emph{the topological entropy of the flow
$g_\bullet$ is, by definition, the topological entropy $h_{\mathrm{top}}(g_1)$
of its time-one map.}

The object of this paper is the following theorem.

\begin{theorem}\label{thm:main}
Assume that $X$ has the following property $(\mathrm{H})$:
\begin{quote}\itshape
whenever $f\in C^\infty(M)$ and $u\in\Dp(M)$ satisfy
$\Lx u=f$ in the sense of distributions, then $u$ is a smooth function.
\end{quote}
Then $h_{\mathrm{top}}(g_1)=0$; equivalently, the flow $g_\bullet$ has zero
topological entropy.
\end{theorem}

Theorem \ref{thm:main} is the contrapositive of a classical and deep
theorem of A.\,G.~Kushnirenko \cite{Kus} (see also the discussion of the
Liv\v{s}ic cohomological equation in \cite{dlLMM}): \emph{a smooth flow
with positive topological entropy admits smooth data $f$ for which the
cohomological equation $\Lx u=f$ has distributional solutions that are
not smooth.} The purpose of this paper is twofold. First, we give a
complete and essentially self-contained derivation of the strong
structural consequences of hypothesis $(\mathrm{H})$; these reduce the
situation to the one in which the only available contradiction is
exactly Kushnirenko's theorem. Second, we record precisely where the
deep input is used, and we prove Kushnirenko's input in the important
special cases in which it can be obtained by soft means
(Appendix~\ref{app:kush}).

Our structural results are summarized in the following theorem, proved
in Sections~\ref{sec:structure} and \ref{sec:entropy}.

\begin{theorem}\label{thm:structure}
Assume $(\mathrm{H})$. Then:
\begin{enumerate}
\item[(a)] every $g_\bullet$-invariant probability measure is of the form
$\mu=\rho\,\omega_0$, where $\omega_0$ is any fixed smooth positive
density on $M$ and $\rho\in C^\infty(M)$, $\rho\ge 0$, satisfies
$\operatorname{div}_{\omega_0}(\rho X)\equiv 0$;
\item[(b)] the flow $g_\bullet$ is minimal;
\item[(c)] there is exactly one $g_\bullet$-invariant probability measure
$\mu_0=\rho_0\omega_0$; it is ergodic, and $\rho_0>0$ everywhere on $M$;
\item[(d)] the space of invariant distributions is the one-dimensional
space $\C\,\mu_0$; moreover $X$ has no zeros and (if $\dim M\ge 2$) the
flow has no periodic orbits;
\item[(e)] $h_{\mathrm{top}}(g_1)=h_{\mu_0}(g_1)$, the metric entropy of
the unique invariant measure.
\end{enumerate}
\end{theorem}

With Theorem \ref{thm:structure} in hand, the proof of Theorem
\ref{thm:main} is completed in Section \ref{sec:kush} by applying
Kushnirenko's theorem to $\mu_0$.

\paragraph{Historical remarks.}
The regularity theory of the cohomological equation
$\Lx u=f$ originates in the work of Liv\v{s}ic, Guivarc'h and others; for
smooth hyperbolic systems the solvability and regularity of solutions is
governed by obstructions over periodic orbits, and one cannot expect
smooth solvability in positive entropy. The precise statement that
positive entropy forces the existence of non-smooth distributional
solutions is due to Kushnirenko \cite{Kus}; it is quoted, and its role
in the regularity theory of the Liv\v{s}ic equation is explained, in the
survey parts of \cite{dlLMM}. Related work on entropy of smooth flows
which we use or cite below includes the variational principle for flows
\cite{Good}, as used in \cite{LWW}; Pesin's entropy formula for flows
\cite{HPT}; and Ruelle-type inequalities for divergence-free fields
\cite{BCF}. Smooth zero-entropy flows with interesting statistical
behavior are constructed in \cite{Dolgopyat}.

\section{Conventions and preliminaries}\label{sec:prelim}

Throughout the paper, $M$ is a compact connected smooth manifold without
boundary and $X$ a $C^\infty$ vector field with flow $g_t$. If
$\dim M=0$ then $M$ is a point, the flow is trivial and
$h_{\mathrm{top}}(g_1)=0$; henceforth $\dim M\ge 1$. If $\dim M=1$, then
$M\cong S^1$; under $(\mathrm{H})$ the field $X$ has no zeros
(Corollary~\ref{cor:invis} below), and a nonvanishing field on $S^1$ has
all orbits periodic with $M$ a single orbit, so $g_1$ is conjugate to a
rotation and $h_{\mathrm{top}}(g_1)=0$. Hence we may and do assume
$\dim M\ge 2$ from Section~\ref{sec:structure} on; all arguments below
are in any case uniform.

\begin{convention}[Distributions and smooth densities]\label{conv:distr}
We fix once and for all a smooth positive density $\omega_0$ on $M$ (for
instance the Riemannian density of some auxiliary metric; no orientability
is required). Since $M$ is compact, the space of test functions is
$C^\infty(M)$ with its Fr\'echet topology, and $\Dp(M)$ denotes its
topological dual. For $h\in C^\infty(M)$ we let
\[
[h]\in \Dp(M),\qquad [h](\varphi):=\int_M h\,\varphi\,d\omega_0 ,
\]
and we say that a distribution $u$ is \emph{smooth} if $u=[h]$ for some
$h\in C^\infty(M)$. This notion does not depend on the choice of
$\omega_0$: if $\omega_1$ is another smooth positive density, then
$[h]_{\omega_0}=[k]_{\omega_1}$ holds if and only if
$k=h\,d\omega_0/d\omega_1$, and $d\omega_0/d\omega_1\in C^\infty(M)$ is
nowhere zero. A (finite complex) measure $\mu$ on $M$ defines the
distribution $T_\mu(\varphi):=\int\varphi\,d\mu$; note that $T_\mu$ is
smooth if and only if $\mu$ has a $C^\infty$ density with respect to
$\omega_0$.
\end{convention}

For a vector field $Y$ we write $\operatorname{div}_{\omega_0}Y$ for the
smooth function determined by
$\mathcal L_Y\omega_0=(\operatorname{div}_{\omega_0}Y)\,\omega_0$; since
$M$ is closed, $\int_M \operatorname{div}_{\omega_0}(Y)\,d\omega_0=0$ for
every smooth field $Y$.

\begin{definition}[Lie derivative of distributions]\label{def:Lieder}
For $u\in\Dp(M)$ define $\Lx u\in\Dp(M)$ by
\begin{equation}\label{eq:Lieder}
\langle \Lx u,\varphi\rangle:=-\langle u,X\varphi\rangle ,
\qquad \varphi\in C^\infty(M).
\end{equation}
Equivalently, $\Lx$ is the transpose of the first-order operator
$\varphi\mapsto -X\varphi$; it is the unique continuous extension of the
Lie derivative acting on smooth densities.
\end{definition}

\begin{lemma}\label{lem:onsmooth}
For $h\in C^\infty(M)$,
\[
\Lx[h]=\bigl[\,Xh+h\operatorname{div}_{\omega_0}X\,\bigr].
\]
In particular, if $X$ is divergence-free with respect to $\omega_0$, then
$\Lx[h]=[Xh]$.
\end{lemma}

\begin{proof}
For $\varphi\in C^\infty(M)$, apply the divergence theorem to the field
$h\varphi X$:
\[
0=\int_M \operatorname{div}_{\omega_0}(h\varphi X)\,d\omega_0
 =\int_M\bigl(\varphi\,Xh+h\,X\varphi
+h\,\varphi\operatorname{div}_{\omega_0}X\bigr)\,d\omega_0 .
\]
Hence $\langle \Lx[h],\varphi\rangle=-\int h\,X\varphi\,d\omega_0
=\int \varphi\bigl(Xh+h\operatorname{div}_{\omega_0}X\bigr)\,d\omega_0$.
\end{proof}

\begin{lemma}\label{lem:flowaction}
The flow acts continuously on $\Dp(M)$ by
$(T_tu)(\varphi):=u(\varphi\circ g_t)$, and
$\frac{d}{dt}T_tu=-T_t\,\Lx u$. In particular, if $\Lx u=0$ then
$T_tu=u$ for all $t$, i.e.\ $u(\varphi\circ g_t)=u(\varphi)$ for all
$\varphi$ and all $t$.
\end{lemma}

\begin{proof}
Since $X$ is smooth and $M$ compact, $t\mapsto \varphi\circ g_t$ is a
$C^\infty$ curve in the Fr\'echet space $C^\infty(M)$, so $t\mapsto T_tu$
is differentiable in $\Dp(M)$ with
$\frac d{dt}\langle T_tu,\varphi\rangle
=u\bigl((X\varphi)\circ g_t\bigr)=\langle T_tu,X\varphi\rangle
=-\langle \Lx (T_tu),\varphi\rangle$.
\end{proof}

\begin{lemma}\label{lem:invmeasure}
Let $\mu$ be a $g_\bullet$-invariant probability measure, i.e.\
$(g_t)_*\mu=\mu$ for all $t$. Then $\Lx T_\mu=0$.
\end{lemma}

\begin{proof}
For each $\varphi$, the function $t\mapsto\int\varphi\circ g_t\,d\mu$
is constant by invariance; differentiating at $t=0$ gives
$\int X\varphi\,d\mu=0$, i.e.\
$\langle \Lx T_\mu,\varphi\rangle=-\int X\varphi\,d\mu=0$.
\end{proof}

\begin{hypothesis}[(H)]\label{hyp:H}
For every $f\in C^\infty(M)$ and every $u\in\Dp(M)$ with
$\Lx u=[f]$, the distribution $u$ is smooth in the sense of Convention
\ref{conv:distr}.
\end{hypothesis}

Taking $f\equiv 0$ in Hypothesis \ref{hyp:H} we obtain the standing
consequence:

\begin{corollary}\label{cor:invdist}
Under $(\mathrm{H})$, every distribution $u$ with $\Lx u=0$ is smooth. In
particular, by Lemma \ref{lem:invmeasure}, for every $g_\bullet$-invariant
probability $\mu$ there exists $\rho\in C^\infty(M)$ with
$T_\mu=[\rho]$, i.e.\ $\mu=\rho\,\omega_0$.
\end{corollary}

\section{Structural consequences of Hypothesis (H)}\label{sec:structure}

In this section we assume Hypothesis \ref{hyp:H}. Recall $\dim M\ge2$.

\begin{lemma}\label{lem:density}
Let $\mu=\rho\,\omega_0$ be a $g_\bullet$-invariant probability with
$\rho\in C^\infty(M)$ (which exists by Corollary \ref{cor:invdist}).
Then:
\begin{enumerate}
\item[(i)] $\rho\ge 0$ and $\operatorname{div}_{\omega_0}(\rho X)\equiv0$,
i.e.\ $X\rho=-\rho\operatorname{div}_{\omega_0}X$;
\item[(ii)] for all $x\in M$ and $t\in\R$,
\begin{equation}\label{eq:transport}
\rho(g_tx)=\rho(x)\exp\Bigl(-\int_0^t
(\operatorname{div}_{\omega_0}X)(g_sx)\,ds\Bigr);
\end{equation}
\item[(iii)] the set $U_\mu:=\{\rho>0\}$ is nonempty, open and
$g_\bullet$-invariant, and $\rho\equiv 0$ on $M\setminus U_\mu$.
\end{enumerate}
\end{lemma}

\begin{proof}
(i) If $\varphi\ge0$ then $\int\rho\varphi\,d\omega_0=\int\varphi\,d\mu\ge0$;
hence the smooth function $\rho$ is nonnegative (otherwise a small
nonnegative test function supported where $\rho<0$ would give a negative
value). By Lemmas \ref{lem:onsmooth} and \ref{lem:invmeasure},
$0=\Lx T_\mu=\Lx[\rho]=[X\rho+\rho\operatorname{div}_{\omega_0}X]$, so
$X\rho=-\rho\operatorname{div}_{\omega_0}X$ pointwise.

(ii) Differentiate $t\mapsto\rho(g_tx)$ using (i):
$\frac d{dt}\rho(g_tx)=(X\rho)(g_tx)
=-\rho(g_tx)(\operatorname{div}_{\omega_0}X)(g_tx)$; solving this linear
equation gives \eqref{eq:transport}.

(iii) $U_\mu$ is open because $\rho$ is continuous, nonempty because
$\mu\neq0$, and invariant by \eqref{eq:transport} (which shows
$\rho(g_tx)=0\iff\rho(x)=0$). The last claim is trivial.
\end{proof}

\begin{lemma}[Krylov--Bogolyubov]\label{lem:KB}
Let $K\subseteq M$ be nonempty, compact and $g_\bullet$-invariant
(i.e.\ $g_t(K)=K$ for all $t$). Then there exists a
$g_\bullet$-invariant probability measure supported in $K$. The same
holds for $K=M$.
\end{lemma}

\begin{proof}
Pick $\nu_0\in\Prob(K)$ and set
$\nu_T:=\frac1T\int_0^T(g_t)_*\nu_0\,dt$. Since $\Prob(M)$ is weakly-$*$
compact ( $M$ compact metric), a subnet $\nu_{T_j}$ converges to some
$\nu\in\Prob(M)$, clearly supported in $K$. For fixed $s$ and test
$\varphi$,
\[
\bigl|\langle (g_s)_*\nu_T-\nu_T,\varphi\rangle\bigr|
=\frac1T\Bigl|\int_T^{T+s}-\int_0^{s}\langle\nu_0,\varphi\circ g_r\rangle\,dr\Bigr|
\le \frac{2|s|\,\|\varphi\|_\infty}{T}\to 0,
\]
so $\nu$ is $g_s$-invariant for each $s$.
\end{proof}

\begin{proposition}\label{prop:opensets}
Assume $(\mathrm{H})$. Every nonempty closed $g_\bullet$-invariant set
$K\subseteq M$ contains a nonempty open $g_\bullet$-invariant subset.
\end{proposition}

\begin{proof}
By Lemma \ref{lem:KB} choose a $g_\bullet$-invariant probability $\nu$
supported in $K$; by Corollary \ref{cor:invdist} and Lemma
\ref{lem:density}, $\nu=\rho\,\omega_0$ with $U:=\{\rho>0\}$ nonempty,
open, invariant. Since $\nu(M\setminus K)=0$ and $M\setminus K$ is open,
we get $\int_{M\setminus K}\rho\,d\omega_0=0$; as $\rho\ge0$ is smooth,
$\rho\equiv0$ on $M\setminus K$, whence $U\subseteq K$.
\end{proof}

\begin{theorem}[Minimality]\label{thm:minimal}
Assume $(\mathrm{H})$. Then the flow $g_\bullet$ is minimal: the only
nonempty closed $g_\bullet$-invariant subset of $M$ is $M$ itself.
Consequently $X$ has no zeros and $g_\bullet$ has no periodic orbits.
\end{theorem}

\begin{proof}
\emph{Step 1: minimal sets exist.} A \emph{minimal set} is a nonempty
compact invariant set containing no proper nonempty compact invariant
subset. The collection of such sets is nonempty by Zorn's lemma: in a
chain ordered by inclusion, the intersection is nonempty (compact sets
with the finite intersection property), compact and invariant.

\emph{Step 2: every minimal set is open.} Let $K$ be minimal. By
Proposition \ref{prop:opensets} there is a nonempty open invariant
$U\subseteq K$. If $x\in U$ then the orbit of $x$ lies in $U$ and is
dense in $K$ (its closure is a nonempty closed invariant subset of the
minimal set $K$), hence $\overline U=K$. Therefore
$\partial U=K\setminus U$ is a closed invariant set; if it were nonempty
it would be a proper nonempty closed invariant subset of $K$ (proper
because $U\neq\emptyset$), contradicting minimality. Hence $K=U$ is
open.

\emph{Step 3: the minimal sets tile $M$.} Distinct minimal sets are
disjoint (their intersection would be a nonempty compact invariant set
contained in both). Being disjoint open nonempty sets, we may assign to
each minimal set a basic open set contained in it; the assignment is
injective, so ( $M$ is second countable) the family
$\{K_j\}_{j\in J}$ of minimal sets is at most countable. The leftover
$L:=M\setminus\bigcup_jK_j$ is closed and invariant. If $L\neq\emptyset$
it contains a minimal set $K$, which must equal some $K_j$; but then
$K_j\subseteq L$ while $K_j\cap\bigcup_iK_i=\emptyset$, a contradiction.
Hence $L=\emptyset$ and $M=\bigsqcup_j K_j$.

\emph{Step 4: conclusion.} Since $M$ is connected, the disjoint open
cover $\{K_j\}$ has exactly one member; thus $M$ is minimal. Finally, a
zero of $X$, or a periodic orbit, would be a nonempty closed invariant
subset of $M$, hence equal to $M$; the first is impossible since
$\dim M\ge2$ (a point is not a manifold of dimension $\ge2$), the second
because an embedded circle has empty interior and cannot equal $M$.
\end{proof}

\begin{theorem}[Unique ergodicity]\label{thm:unique}
Assume $(\mathrm{H})$. Then there exists exactly one
$g_\bullet$-invariant probability measure $\mu_0$. Writing
$\mu_0=\rho_0\,\omega_0$ with $\rho_0\in C^\infty(M)$, one has
$\rho_0>0$ everywhere; $\mu_0$ is ergodic; and every finite
$g_\bullet$-invariant complex measure is a complex multiple of $\mu_0$.
\end{theorem}

\begin{proof}
\emph{Existence:} Lemma \ref{lem:KB} with $K=M$.

\emph{Uniqueness:} Let $\mu_1=\rho_1\omega_0$ and $\mu_2=\rho_2\omega_0$
be two invariant probabilities (densities as in Lemma \ref{lem:density}).
Put $\rho:=\rho_1+\rho_2$, $U:=\{\rho>0\}$, and
$\kappa:=\rho_1/\rho\in C^\infty(U,[0,1])$. By \eqref{eq:transport},
applied to $\rho_1$ and $\rho_2$ (which involve the \emph{same}
cocycle $\int_0^t\operatorname{div}_{\omega_0}X\,ds$),
\[
\kappa(g_tx)=\frac{\rho_1(g_tx)}{\rho_1(g_tx)+\rho_2(g_tx)}
=\frac{\rho_1(x)}{\rho_1(x)+\rho_2(x)}=\kappa(x)
\qquad (x\in U).
\]
Fix $x\in U$. By Theorem \ref{thm:minimal} the orbit of $x$ is dense in
$M$ and it stays in $U$. Given $y\in U$, choose $t_j$ with
$g_{t_j}x\to y$; continuity of $\kappa$ on the open set $U$ gives
$\kappa(y)=\lim_j\kappa(g_{t_j}x)=\kappa(x)$. Hence $\kappa\equiv c$ on
$U$ for some constant $c$, i.e.\ $\rho_1\equiv c\rho$ and
$\rho_2\equiv(1-c)\rho$ on $U$, and all three densities vanish outside
$U$. If $c=0$ then $\rho_1\equiv0$, so $\mu_1=0$, a contradiction; if
$c=1$ then $\rho_2\equiv0$; if $0<c<1$ then
$\mu_1=\frac{c}{1-c}\,\mu_2$ and since both are probabilities,
$\mu_1=\mu_2$. Hence $\mu_1=\mu_2$.

\emph{Full support:} $F:=M\setminus\{\rho_0>0\}$ is closed and
invariant (Lemma \ref{lem:density}(iii)); if $F\neq\emptyset$, Lemma
\ref{lem:KB} gives an invariant probability supported in $F$, which by
uniqueness equals $\mu_0$; but
$\mu_0(F)\le\int_F\rho_0\,d\omega_0=0$ and yet that measure has total
mass $1$ --- contradiction. Hence $\rho_0>0$ everywhere.

\emph{Ergodicity:} if $A$ is invariant Borel with
$0<\mu_0(A)<1$, then $\mu_0(\cdot\cap A)/\mu_0(A)$ is an invariant
probability distinct from $\mu_0$, contradicting uniqueness.

\emph{Invariant measures:} a finite positive invariant measure
$\sigma\neq0$ normalizes to an invariant probability, hence equals
$\sigma(M)\mu_0$; a finite complex invariant measure decomposes by Jordan
into a combination of finite positive invariant measures, hence is a
complex multiple of $\mu_0$.
\end{proof}

\begin{corollary}[Invariant distributions]\label{cor:invdim}
Assume $(\mathrm{H})$. The space of invariant distributions
$\{u\in\Dp(M): \Lx u=0\}$ is the one-dimensional space
$\C\,T_{\mu_0}=\{[c\rho_0]:c\in\C\}$. Consequently $X$ has no zeros and
$g_\bullet$ has no periodic orbits.
\end{corollary}

\begin{proof}
If $\Lx u=0$, then $u=[h]$ for some $h\in C^\infty(M)$ by
Corollary~\ref{cor:invdist}, and $T_tu=u$ for all $t$ by Lemma
\ref{lem:flowaction}; unwinding the definitions,
$(g_t)_*(h\omega_0)=h\omega_0$ for all $t$, i.e.\ the complex measure
$h\,\omega_0$ is invariant. By Theorem \ref{thm:unique},
$h\,\omega_0=c\,\mu_0=c\,\rho_0\omega_0$, so $u=[c\rho_0]$.

If $p$ were a zero of $X$, then $T_t\delta_p=\delta_p$ and
$\Lx\delta_p=0$ by Lemma \ref{lem:flowaction}, so
$\delta_p=[c\rho_0]$; as measures this says $\delta_p=c\rho_0\omega_0$,
impossible because $\delta_p(\{p\})=1$ while any smooth density gives
mass $0$ to a point. Similarly, the normalized arc measure along a
periodic orbit $\gamma$ is $g_\bullet$-invariant, hence by the above
equals $\mu_0$; but it is supported on the embedded circle $\gamma$,
while $\mu_0(\gamma)=\int_\gamma\rho_0\,d\omega_0=0$ when
$\dim M\ge2$ --- contradiction.
\end{proof}

\section{Entropy under Hypothesis (H)}\label{sec:entropy}

For a homeomorphism $T$ of a compact metric space and a $T$-invariant
probability $\nu$, let $h_\nu(T)$ denote the Kolmogorov--Sinai entropy.
As recalled in the introduction, the topological entropy of the flow
$g_\bullet$ is, by definition, $h_{\mathrm{top}}(g_1)$ \cite{LWW}. We
record the variational principle for flows in the form in which we use
it; it is stated (with proof, based on Goodman's theorem \cite{Good} and
Abramov's formula) for instance in \cite{LWW}, where it appears as
follows: \emph{$h_{\mathrm{top}}(g_1)$ coincides with the supremum of
$h_m(g_1)$ over all $g_\bullet$-invariant ergodic probability measures
$m$.} (Note that a $g_\bullet$-invariant measure is in particular
$g_1$-invariant, since $g_1$ commutes with the flow.)

\begin{lemma}[Variational principle for flows; \cite{Good,LWW}]\label{lem:VP}
$h_{\mathrm{top}}(g_1)=\sup\{h_m(g_1):\, m$ a $g_\bullet$-invariant
ergodic probability$\}$.
\end{lemma}

\begin{theorem}\label{thm:entropy}
Assume $(\mathrm{H})$, and let $\mu_0=\rho_0\omega_0$ be the unique
invariant probability of Theorem \emph{\ref{thm:unique}}. Then
\[
h_{\mathrm{top}}(g_1)=h_{\mu_0}(g_1).
\]
\end{theorem}

\begin{proof}
By Theorem \ref{thm:unique}, $\mu_0$ is the unique
$g_\bullet$-invariant probability and is ergodic; hence the supremum in
Lemma \ref{lem:VP} is taken over the singleton $\{\mu_0\}$.
\end{proof}

\begin{remark}\label{rem:pesin}
Suppose $(\mathrm{H})$ holds and $h_{\mathrm{top}}(g_1)>0$. Then
$h_{\mu_0}(g_1)>0$ with $\mu_0$ a $C^\infty$ invariant measure, and
Ruelle's inequality / Pesin's entropy formula for $C^{1+\alpha}$
divergence-free (more generally, absolutely continuous invariant
measure) flows applies: $h_{\mu_0}(g_1)$ equals the integral of the sum
of positive Lyapunov exponents \cite{HPT,BCF}. Thus $\mu_0$ has a
positive Lyapunov exponent and Pesin unstable manifolds exist almost
everywhere. As explained in Appendix \ref{app:kush}, even this structure
does not contradict $(\mathrm{H})$ by soft means, and the decisive input
remains Kushnirenko's theorem.
\end{remark}

\section{Kushnirenko's theorem and the proof of Theorem
\ref{thm:main}}\label{sec:kush}

The following is the classical theorem of A.\,G.~Kushnirenko
\cite{Kus}; its role in the regularity theory of the cohomological
equation is explained in \cite{dlLMM}. We formulate it with the
conventions of Section \ref{sec:prelim}.

\begin{theorem}[Kushnirenko]\label{thm:kush}
Let $N$ be a compact smooth manifold, $Y$ a $C^\infty$ vector field on
$N$ with flow $\Phi_\bullet$, and suppose some
$\Phi_\bullet$-invariant probability measure $\mu$ has positive metric
entropy $h_\mu(\Phi_1)>0$. Then there exist a function $f\in C^\infty(N)$
and a distribution $u\in\Dp(N)$ such that $\mathcal L_Yu=[f]$ and such
that $u$ is not smooth (i.e.\ $u$ is not of the form $[h]$,
$h\in C^\infty(N)$).
\end{theorem}

This is a deep result. Its proof in \cite{Kus} is based on an
$\varepsilon$-entropy (finite-dimensional approximation) argument:
positive metric entropy forces exponential growth of the orbit
complexity of smooth functions in the Sobolev scale, and this in turn
produces smooth data whose distributional primitives exist but are not
smooth. In Appendix \ref{app:kush} we prove Theorem \ref{thm:kush} in
the special cases in which soft methods suffice (existence of periodic
orbits or of invariant measures which are not smooth densities, and the
case of hyperbolic invariant measures via Katok's closing lemma), and we
explain why, under Hypothesis $(\mathrm{H})$, precisely the remaining
case is needed.

\begin{proof}[Proof of Theorem \emph{\ref{thm:main}}]
Assume $(\mathrm{H})$ and suppose, for contradiction, that
$h_{\mathrm{top}}(g_1)>0$. By Theorem \ref{thm:entropy},
$h_{\mu_0}(g_1)=h_{\mathrm{top}}(g_1)>0$, where $\mu_0$ is the unique
$g_\bullet$-invariant probability. Applying Theorem \ref{thm:kush} with
$N=M$, $Y=X$, $\mu=\mu_0$, we obtain $f\in C^\infty(M)$ and a
non-smooth $u\in\Dp(M)$ with $\Lx u=[f]$. This contradicts Hypothesis
\ref{hyp:H}. Hence $h_{\mathrm{top}}(g_1)=0$.
\end{proof}

\begin{remark}
Combining Theorems \ref{thm:structure} and \ref{thm:main}: under
$(\mathrm{H})$ the flow is minimal, uniquely ergodic, preserves a smooth
volume, and has zero topological (hence, by Lemma \ref{lem:VP}, zero
metric) entropy. Conversely, the reduction shows that in order to
contradict $(\mathrm{H})$ in positive entropy one may assume, without
loss of generality, that all of these rigidity properties hold; the only
available contradiction is then Kushnirenko's theorem. Note that the
class of minimal smooth flows with smooth invariant volume and positive
topological entropy is not vacuous: suspensions of Herman's celebrated
minimal, volume-preserving, positive-entropy diffeomorphisms (see
\cite{LWW} and the references therein to Herman's construction) have
exactly these properties, although by Theorem \ref{thm:main} they fail
$(\mathrm{H})$.
\end{remark}

\section{Appendix: remarks on Kushnirenko's theorem}\label{app:kush}

We collect here (i) soft proofs of Theorem \ref{thm:kush} in special
cases, and (ii) an explanation of why the general case is genuinely
needed under Hypothesis $(\mathrm{H})$.

\paragraph{(a) The case of a non-smooth invariant measure.}
Suppose $N$, $Y$, $\Phi_\bullet$ are as in Theorem \ref{thm:kush} and
that \emph{some} $\Phi_\bullet$-invariant probability $\nu$ is not of
the form $\rho\,\omega_N$ with $\rho\in C^\infty(N)$ ( $\omega_N$ any
smooth positive density). Then $T_\nu$ is an invariant distribution
(Lemma \ref{lem:invmeasure}) which is not smooth (Convention
\ref{conv:distr}), and the conclusion of Theorem \ref{thm:kush} holds
with $f\equiv 0$ and $u=T_\nu$. In particular this applies if
$\Phi_\bullet$ has a fixed point (take $\nu=\delta_p$) or a periodic
orbit whose normalized arc measure is not a smooth density (always the
case when $\dim N\ge2$), cf.\ Corollary \ref{cor:invdim}.

\paragraph{(b) The case of a hyperbolic invariant measure.}
Suppose $\mu$ is an ergodic $\Phi_\bullet$-invariant probability with
$h_\mu(\Phi_1)>0$ all of whose Lyapunov exponents, off the flow
direction, are nonzero. By the closing lemma of Katok for flows
\cite{Katok,KH}, $\Phi_\bullet$ then possesses hyperbolic periodic
orbits; their arc measures are invariant and singular with respect to
any smooth density ( $\dim N\ge2$), so case (a) applies. Thus
Theorem \ref{thm:kush} holds in this case as well. (For the entropy
formula $h_\mu(\Phi_1)=\int\sum\lambda^+\,d\mu$ for flows with smooth
invariant measures, which guarantees a positive exponent once the
entropy is positive, see \cite{HPT,BCF}.)

\paragraph{(c) Why the general case is needed.}
Under Hypothesis $(\mathrm{H})$, routes (a) and (b) are closed:
by Corollary \ref{cor:invdim} there are no fixed points and no periodic
orbits, and by Corollary \ref{cor:invdist} \emph{every} invariant
probability has a smooth density. Moreover minimal, smooth,
volume-preserving flows with positive topological entropy exist
(suspensions of Herman's examples \cite{LWW}), so no soft mechanism can
produce a contradiction. The general case of Theorem
\ref{thm:kush}---where the flow may be minimal and uniquely ergodic
with smooth invariant volume---is exactly Kushnirenko's contribution in
\cite{Kus}, and we have cited it as such rather than reproduce its
(considerable) $\varepsilon$-entropy argument. We emphasize that, by
Theorem \ref{thm:structure}, in the proof of Theorem \ref{thm:main} this
input is applied only in the reduced setting of a minimal, uniquely
ergodic flow with smooth invariant volume and positive entropy.

\paragraph{(d) Indication of the mechanism in \cite{Kus}.}
For the reader's orientation (this is a second-hand description, not a
proof): the argument of \cite{Kus} shows that if all distributional
solutions of $\mathcal L_Yu=f$ with smooth $f$ were smooth, then the
orbit complexity of smooth functions under the flow, measured by
$\varepsilon$-entropies of finite-dimensional spectral projections,
could grow at most subexponentially; Kushnirenko's finite-dimensional
approximation estimate then bounds the metric entropy of every
invariant measure by zero. The contrapositive is Theorem
\ref{thm:kush}. Full details are in \cite{Kus}; the regularity-theoretic
context appears in \cite{dlLMM}.

\begin{thebibliography}{99}

\bibitem{Kus} A.\,G.~Kushnirenko, \emph{An estimate of the entropy of
smooth dynamical systems}, Funkcional. Anal. i Prilo\v{z}en. \textbf{8}
(1974); English transl., Functional Analysis and Its Applications
\textbf{8} (1974). (Source of Theorem \ref{thm:kush}.)

\bibitem{dlLMM} R.~de la Llave, J.\,M.~Marco, R.~Moriy\'on,
\emph{Canonical perturbation theory of Anosov systems, and regularity
results for the Liv\v{s}ic cohomology equation}, Ann. of Math. (2)
\textbf{123} (1986), 537--611.

\bibitem{Good} T.\,N.\,T.~Goodman, \emph{Relating topological entropy and
measure entropy}, Bull. London Math. Soc. \textbf{3} (1971), 176--180.

\bibitem{Katok} A.~Katok, \emph{Lyapunov exponents, entropy and periodic
orbits for diffeomorphisms}, Inst. Hautes \'Etudes Sci. Publ. Math.
\textbf{51} (1980), 137--173.

\bibitem{KH} A.~Katok, B.~Hasselblatt, \emph{Introduction to the Modern
Theory of Dynamical Systems}, Cambridge University Press, 1995.

\bibitem{Walters} P.~Walters, \emph{An Introduction to Ergodic Theory},
Graduate Texts in Math. \textbf{79}, Springer, 1982.

\bibitem{LWW} \emph{Topological entropies of equivalent smooth flows}
(on the definition of flow entropy as the entropy of the time-one map,
the variational principle for flows over flow-invariant ergodic
measures, and Herman's minimal positive-entropy examples),
\url{https://exa.ai/library/publication/lbqpyt46khy}

\bibitem{HPT} \emph{On the entropy of conservative flows} (Pesin's
entropy formula for flows),
\url{https://ar5iv.labs.arxiv.org/html/1002.2366}

\bibitem{BCF} \emph{Lyapunov exponents and entropy for divergence-free
Lipschitz vector fields} (Ruelle-type inequality),
\url{https://link.springer.com/content/pdf/10.1007/s40879-023-00611-6.pdf}

\bibitem{Dolgopyat} D.~Dolgopyat, \emph{Smooth zero entropy flows
satisfying the classical central limit theorem},
J. \'Ec. polytech. Math., \url{https://jep.centre-mersenne.org/articles/10.5802/jep.283/}

\bibitem{LLE} \emph{Entropy and periodic orbits for equivalent smooth
flows}, \url{https://ar5iv.labs.arxiv.org/html/1111.0111}

\end{thebibliography}

\end{document}
```

**Summary of the argument.**

1. **Conventions.** Fixing a smooth density $\omega_0$, distributions are identified with generalized densities, and $\mathcal L_Xu$ is defined by $\langle\mathcal L_Xu,\varphi\rangle=-\langle u,X\varphi\rangle$; under this (standard) convention, smooth densities satisfy $\mathcal L_X[h]=[Xh+h\,\mathrm{div}_{\omega_0}X]$ and every flow-invariant measure defines an invariant distribution (Lemmas 2.1–2.3).

2. **Structural consequences of (H)** (fully proved, essentially self-contained): invariant measures have $C^\infty$ densities satisfying the invariance PDE, whose positivity sets are open and invariant; every nonempty closed invariant set therefore contains a nonempty open invariant subset; minimal sets are open, and connectedness forces minimality; a first-integral argument on ratios of densities then gives unique ergodicity with an everywhere-positive smooth density, ergodicity, a one-dimensional space of invariant distributions, and the absence of fixed points and periodic orbits (Theorems 3.4–3.6).

3. **Entropy** (Section 4): by the variational principle for flows (Goodman, as stated in [exa.ai]), $h_{\mathrm{top}}(g_1)=h_{\mu_0}(g_1)$, the metric entropy of the unique smooth invariant measure.

4. **Conclusion** (Section 5): if this were positive, Kushnirenko's classical theorem [Kus74] (positive entropy forces smooth data $f$ with non-smooth distributional solutions of $\mathcal L_Xu=f$; soft special cases proved in the Appendix, including the hyperbolic-measure case via Katok's closing lemma, with Pesin/Ruelle input cited from [ar5iv:1002.2366] and [link.springer.com]) would contradict hypothesis (H). Hence the entropy vanishes.